\documentclass[12pt]{article}
\usepackage[T2A,T1]{fontenc}
\usepackage[utf8]{inputenc}
\usepackage[english,ukrainian]{babel}
\usepackage[left=30mm,right=15mm,top=20mm,bottom=20mm]{geometry}
\usepackage{setspace}
\usepackage{parskip}
\usepackage{tikz}
\usetikzlibrary{svg.path}

\onehalfspacing
\setlength{\parindent}{1.25cm}

\newcommand{\tryzub}[1][1]{%
  \begin{tikzpicture}[scale=0.03*#1, yscale=-1, baseline=(current bounding box.center)]
    % Shield: transparent background, black stroke
    \filldraw[fill=white, draw=black, line width=1.5pt]
      svg {m5 5h650v689c0 48-29 97-76 117l-251 105-251-105c-44-20-76-65-72-117z};

    % Right half of Tryzub (Black)
    \fill[color=black] 
      svg {m329 53c-6 4-2 396 0 401 12 43 29 81 48 112-104 31-63 146-48 287 7-12 17-21 28-28 43-34 70-81 79-132h133v-580c-148 88-132 213-148 361 59-10 75 76-3 83-92-149-59-257-59-419 0-37-9-62-30-85zm200 143v297h-22c-6-23-21-41-42-51 8-85 10-189 64-246zm-22 337h22v120h-89c0-19-3-39-8-58 36-4 68-29 75-62zm-114 71c4 16 6 33 6 49h-50c1-27 19-44 44-49zm-44 89h46c-7 31-23 61-46 85z};

    % Left (mirrored) half of Tryzub (Black)
    \fill[color=black] 
      svg {m331 53c6 4 2 396 0 401-12 43-29 81-48 112 104 31 63 146 48 287-7-12-17-21-28-28-43-34-70-81-79-132h-133v-580c148 88 132 213 148 361-59-10-75 76 3 83 92-149 59-257 59-419 0-37 9-62 30-85zm-200 143v297h22c6-23 21-41 42-51-8-85-10-189-64-246zm22 337h-22v120h89c0-19 3-39 8-58-36-4-68-29-75-62zm114 71c-4 16-6 33-6 49h50c-1-27-19-44-44-49zm44 89h-46c7 31 23 61 46 85z};
  \end{tikzpicture}%
}


\begin{document}

\begin{center}
  \tryzub[0.8]\\[0.3cm]
  {\selectlanguage{ukrainian}\small\MakeUppercase{Державна судова адміністрація України}\\}
  {\selectlanguage{ukrainian}\small\bfseries\MakeUppercase{Державне підприємство «Інформаційні судові системи»}\\}
  \vspace{0.2cm}
  \hrule height 1.5pt
  \vspace{0.5cm}
  {\selectlanguage{ukrainian}\Large\bfseries НАКАЗ\\}
\end{center}

\vspace{0.2cm}

\noindent
\parbox[t]{0.33\textwidth}{\selectlanguage{ukrainian} 09.01.2023}%
\parbox[t]{0.33\textwidth}{\centering\selectlanguage{ukrainian} Київ}%
\parbox[t]{0.33\textwidth}{\raggedleft\selectlanguage{ukrainian} № 3\_\allowbreak{}ОД}

\vspace{0.8cm}

\noindent\textbf{Про затвердження Положення про розподіл відповідальності, прав та обов'язків адміністраторів ЄСІКС}

\vspace{0.5cm}

З метою організаційного забезпечення захисту інформації в Єдиній судовій інформаційно-комунікаційній системі «ЄСІКС» (далі -- ЄСІКС), врегулювання основних аспектів діяльності адміністраторів щодо порядку доступу до інформаційних ресурсів, забезпечення їх взаємодії, порядку оброблення інформації в ЄСІКС, відповідно до Законів України «Про інформацію», «Про захист інформації в інформаційно-комунікаційних системах», «Про захист персональних даних», а також нормативних документів у сфері технічного захисту інформації (НД ТЗІ 2.5-004-99, НД ТЗІ 2.3-025-24, НД ТЗІ 3.7-003-23) та на виконання контролів безпеки AT, IR, MA, PS, PE, PT,

\noindent\textbf{НАКАЗУЮ:}

1. Затвердити Положення про розподіл відповідальності, прав та обов'язків адміністраторів Єдиної судової інформаційно-комунікаційної системи «ЄСІКС» (далі -- Положення), що додається (Додаток 1).

2. Затвердити Політики та інструкції з інформаційної безпеки, що регулюють:
\begin{itemize}
  \item Програму підвищення обізнаності та навчання з безпеки (AT-2, AT-3);
  \item Процедури реагування на інциденти інформаційної безпеки (IR-4, IR-5);
  \item Регламент технічного обслуговування компонентів ЄСІКС (MA-2, MA-4);
  \item Правила кадрової безпеки та нерозголошення (PS-3, PS-4);
  \item Режим фізичного та екологічного захисту приміщень (PE-2, PE-3);
  \item Політики конфіденційності та згоди на обробку персональних даних (PT-2, PT-4).
\end{itemize}

3. Керівникам структурних підрозділів підприємства забезпечити неухильне дотримання вимог затвердженого Положення та контроль за його дотриманням підлеглими працівниками, яким надано право доступу до ЄСІКС.

4. Керівнику СЗІ Дмитро КОВАЛЕНКО забезпечити ознайомлення з вимогами цього Положення працівників, які визначені адміністраторами та користувачами ЄСІКС, під особистий підпис із внесенням відповідних записів до журналу інструктажів.

5. Контроль за виконанням цього наказу залишаю за собою.

\vspace{1.5cm}

\noindent
\parbox[t]{0.5\textwidth}{\selectlanguage{ukrainian}\bfseries В. о. генерального директора}%
\parbox[t]{0.5\textwidth}{\raggedleft\selectlanguage{ukrainian}\bfseries Олексій ГРИЦЕНКО}

\newpage

\begin{flushright}
ДОДАТОК 1\\
до наказу Державне підприємство «Інформаційні судові системи»\\
від 09.01.2023 р. № 3\_\allowbreak{}ОД
\end{flushright}

\begin{center}
  \textbf{\Large\selectlanguage{ukrainian}ПОЛОЖЕННЯ}\\
  \textbf{\large\selectlanguage{ukrainian}про розподіл відповідальності, прав та обов'язків адміністраторів Єдиної судової інформаційно-комунікаційної системи «ЄСІКС»}
\end{center}

\vspace{0.5cm}

\section{Терміни та визначення}
У цьому документі використовуються терміни та визначення згідно з державними нормативними документами в сфері захисту інформації та стандартів НД ТЗІ 1.1-003-99.
\begin{description}
  \item[Адміністратори] -- працівники підприємства, посадові обов'язки яких передбачають здійснення адміністрування або які наділені відповідними повноваженнями.
  \item[Служба захисту інформації (СЗІ)] -- підрозділ, що забезпечує функціонування КСЗІ.
  \item[Інцидент безпеки] -- подія, що призвела або може призвести до порушення конфіденційності, цілісності чи доступності інформації.
\end{description}

\section{Загальні положення}
1.1. Вимоги цього Положення є обов'язковими для виконання всіма адміністраторами та користувачами ЄСІКС.
1.2. Виконання вимог є необхідним для забезпечення безпеки інформації та безперебійного функціонування ЄСІКС.

\section{Ролі адміністраторів та їх обов'язки}
В ЄСІКС призначаються такі адміністратори:
\begin{itemize}
  \item \textbf{Адміністратор домену} -- забезпечує адміністрування Active Directory, керує обліковими записами та групами користувачів.
  \item \textbf{Адміністратор мережі} -- відповідає за керування активним мережевим обладнанням, портами та транзитом трафіку.
  \item \textbf{Адміністратор систем та прикладного ПЗ} -- відповідає за сервери та прикладне програмне забезпечення.
  \item \textbf{Адміністратор баз даних} -- забезпечує стабільну роботу СУБД та цілісність даних.
  \item \textbf{Адміністратор безпеки} -- здійснює контроль за станом безпеки інформації та роботою інших адміністраторів.
  \item \textbf{Адміністратор резервного копіювання} -- відповідає за резервне копіювання та відновлення даних.
\end{itemize}

Всі адміністратори зобов'язані проходити регулярне навчання (AT-3), негайно повідомляти про інциденти (IR-4) та дотримуватись регламентів обслуговування (MA-2).

\end{document}
